\documentclass[11pt]{article}
\usepackage[T1]{fontenc}
\usepackage{lmodern,amsmath,amssymb,amsthm,mathtools,microtype}
\usepackage[margin=1in]{geometry}
\usepackage{booktabs,enumitem,xcolor,fancyhdr}
\usepackage[colorlinks=true,linkcolor=blue!45!black,urlcolor=blue!45!black,citecolor=blue!45!black]{hyperref}
\setlength{\headheight}{14pt}
\pagestyle{fancy}\fancyhf{}\fancyhead[L]{Coherent recurrent communication}
\fancyhead[R]{Astra Pro 1}\fancyfoot[C]{\thepage}
\setlength{\parskip}{5pt}\setlength{\parindent}{0pt}
\setlist{itemsep=3pt,topsep=4pt}
\newtheorem{theorem}{Theorem}[section]
\newtheorem{lemma}[theorem]{Lemma}
\newtheorem{proposition}[theorem]{Proposition}
\newtheorem{corollary}[theorem]{Corollary}
\newcommand{\Tr}{\operatorname{Tr}}
\newcommand{\rank}{\operatorname{rank}}
\newcommand{\id}{\operatorname{id}}
\newcommand{\ket}[1]{\lvert #1\rangle}
\newcommand{\bra}[1]{\langle #1\rvert}
\newcommand{\norm}[1]{\lVert #1\rVert}
\newcommand{\E}{\mathbb E}
\newcommand{\Prob}{\mathbb P}
\newcommand{\Ic}{I_{\mathrm c}}
\newcommand{\Qcap}{\mathsf Q}
\newcommand{\SWAP}{\operatorname{SWAP}}
\title{Spectral clipping for one recurrent quantum bath\\
A nearly quadratic small-angle communication lower bound}
\author{Astra Pro 1\\Source-complete Cycle9 checkpoint}
\date{22 October 2026, Australia/Brisbane}
\begin{document}
\maketitle
\begin{abstract}
The completed nonlinear response theorem permits temporal blocks much longer
than its original positive linear input. We prove that this gap is real for
that input family: its explicit operator has a positive Rayleigh quotient
of order \(k^{3/2}\), and no normalized repair depending only on total input
charge can remove the resulting restriction.

The same operator nevertheless has exponentially small reference-weighted
spectral tails on scale \(k\). An even-endpoint moment count proves this for
all lengths, allowing arbitrary controlled color permutations. Taking the
positive part of a signed affine input and normalizing it changes its exact
operator drift by an exponentially small trace-weighted amount. A global
chi-square entropy bound avoids a small operator-norm expansion.

For every fixed \(0<\beta\le1/16\), with
\(k=\lfloor[16\beta^2\ln(e/\beta^2)]^{-1}\rfloor\), this positive input has
strict exact inward bulk drift and entropy above \(3k/2\). The preserved
root-retaining coding operation gives
\[
 \Qcap(V_\beta,\tau)\ge\frac32+
 \frac{\ln2}{65536}\frac{\beta^2}{\ln(e/\beta^2)}.
\]
One deterministic unassisted code sequence works pointwise for every fixed
normal initializer in full diamond norm. All uses share the same bath.
There is no terminal-return condition, postselection or efficient-code claim.
The logarithmic loss, sharp capacity behavior and feedback equality remain open.
\end{abstract}
\newpage\tableofcontents\newpage

\section{Physical operation, source dependencies and the new theorem}

The physical bath and device remain
\[
 E=\mathbb C^2_{\rm seed}\otimes\mathcal F(\mathbb C^2),\qquad
 U=\begin{pmatrix}R&S_1&S_2\\S_1^*&0&0\\S_2^*&0&0\end{pmatrix},
 \qquad V_\beta=\sqrt{1-\beta^2}\,U+i\beta I.
\]
The operator \(S_j\) prepends color \(j\). The operator \(R\) exchanges
the two vacuum seeds and vanishes off vacuum. Thus \(U=U^*=U^{-1}\).
Let \(L_E\) be bath depth, \(q_i=I-\ket0\bra0\), \(Q=\sum_{i=0}^{k-1}q_i\),
and \(W_k=V_{k-1,E}\cdots V_{0,E}\). The conserved charge is \(L_E+Q\),
including the actual vacuum exchange.

The actual normal initializer and reference are
\[
 \tau_\pi=\frac{I_2}{2}\otimes\frac23\sum_w6^{-|w|}\ket w\bra w,
 \qquad \eta=\operatorname{diag}(3/4,1/8,1/8).
\]
Its depth mass is \((2/3)3^{-h}\), uniform over \(2^{h+1}\) vertices.
Every forward call uses the same evolving infinite bath. No bath access,
reset, fresh bath, preshared entanglement, postselection, backward quantum
wire or free timing resource is supplied. All calls, outcomes and retained
registers count. The coding contract imposes no terminal bath return.

Put
\[
 b=\beta^2,\quad a=1-b,\quad \ell=\ln2,\quad
 \sigma=\operatorname{diag}(1/2,1/4,1/4)^{\otimes k},\quad
 A=Q-k/2.
\]
In particular \(-\log_2\sigma=kI+Q\).
The exact Cycle8 archive is preserved unchanged, including its Cycle7 and
earlier source chain. Its explicit Hermitian input operator \(H=H_k\)
and response theorem are used with their actual proofs, not reconstructed
from a summary. Appendix~\ref{app:response} reproduces the relevant
response proof; the complete preceding native source supplies the all-length
token construction and its sharp variational calculation.

Here are the inherited facts needed in the new argument:
\begin{align}
 [H,Q]&=[H,\sigma]=0,\qquad
 \Tr\sigma H=\Tr\sigma QH=0,\notag\\
 H&=\sum_{i<j}H_{ij},\qquad \norm{H_{ij}}\le2,\label{eq:newendpoint}\\
 v=v_k&=\Tr(\sigma H^2),\qquad
 \frac{k^2}{32}\le v\le\frac{k(k-1)}8.\label{eq:newv}
\end{align}
Every \(H_{ij}\) changes exactly the binary charge bits \(i,j\), exchanging
\(01\) and \(10\); it has no other binary-charge block. It may act on many
spectator color wires. This is not a two-site locality assumption.
Conditional on the full charge pattern its oriented color block is
\(P(I-\SWAP)\), \(P/2\), or zero, according as the unmatched earlier
push count \(h\) exceeds, equals, or is below the middle deficit \(d\).
The original token algorithm fixes the permutation \(P\) exactly.

Define the bounded physical displacement, without subtracting possibly
infinite depth expectations, by
\[
 G_\beta=Q-W_k^*QW_k,\qquad \norm{G_\beta}\le k,\qquad
 D_b(X)=\Tr_A[(X\otimes I)G_\beta].
\]
The partial-trace bound
\begin{equation}\label{eq:driftnorm}
 \norm{D_b(X)}\le k\norm X_1
\end{equation}
holds for every Hermitian input \(X\), on the full infinite bath.
On initial depths at least \(k\), the exact inherited identities are
\begin{align}
 D_b(\sigma)&=0,\qquad D_b(\sigma A)=\frac{ak}{2}I,\notag\\
 D_b(\sigma H)&=baF_k(b),\qquad
 \norm{F_k(b)-vI}\le8b(k-2)v.\label{eq:newresponse}
\end{align}
The last bound is uniform over the entire coherent/entangled bath prefix.
The bulk identities are not root identities.

\begin{theorem}[Positive spectral repair and reliable communication]
\label{thm:newrate}
For every fixed \(0<\beta\le1/16\), set
\begin{equation}\label{eq:newk}
 T=\ln(e/b),\qquad k=\left\lfloor\frac1{16bT}\right\rfloor,\qquad
 m=\frac{3\ell v_k}{32},\qquad \gamma=\frac{\ell}{4}.
\end{equation}
Define, on the finite \(3^k\)-dimensional input space,
\begin{equation}\label{eq:newrho}
 B=I-\gamma bH+\frac{4m}{k}b^2A,\qquad
 \epsilon=\Tr(\sigma B_-),\qquad
 \rho=\frac{\sigma B_+}{1+\epsilon},
\end{equation}
where \(B_+=\max(B,0)\) and \(B_-=\max(-B,0)\).
Then \(\rho\) is a normalized positive charge-preserving density and
\begin{align}
 S(\rho)-3k/2&\ge\frac{\ell v_k b^2}{64},\label{eq:newentropy}\\
 P_{\ge k}D_b(\rho)P_{\ge k}
 &\le-\kappa P_{\ge k},\qquad
 \kappa=\frac{\ell v_k b^2(1-b)}{128}>0.\label{eq:exactdrift}
\end{align}
Consequently one deterministic unassisted encoder/decoder sequence,
independent of the initializer, achieves
\begin{equation}\label{eq:newrate}
 \Qcap(V_\beta,\tau)\ge\frac{S(\rho)}k
 \ge\frac32+\frac{\ell v_k\beta^4}{64k}
 \ge\frac32+\frac{\ln2}{65536}
                  \frac{\beta^2}{\ln(e/\beta^2)}
\end{equation}
with pointwise vanishing full-diamond error for every fixed normal bath
initializer \(\tau\). At the actual initializer,
\[
 \lim_{n\to\infty}\frac1{nk}
 \Ic(\rho^{\otimes n},\mathcal N_{nk}^{\tau_\pi})=\frac{S(\rho)}k.
\]
The parameters \(\beta,k,\rho\) are fixed before the coding horizon \(n\).
The channel is the actual same-bath channel, not a tensor power of fresh channels.
\end{theorem}

The positive part in \eqref{eq:newrho} defines an input density. It is not a
probabilistic physical filter or a successful branch selected for credit.
The inherited coding operation constructs an isometric encoder from this
density's spectral types. Its rank is \(d=\rank B_+\le3^k\); clipping
may create zero eigenvalues, which are omitted from those types.

\section{A positive-spectrum obstruction for the old linear family}

Let
\[
 C_k=\sum_{n=0}^{k-2}(k-1-n)\Prob(S_n=0)
     =\frac{4k^{3/2}}{3\sqrt{2\pi}}+O(\sqrt k),
\]
where \(S_n\) is a simple symmetric random walk. This exact collision
quantity and asymptotic were proved in Cycle7. In that source
\(v_k=k(k-1)/8-C_k/8\).

\begin{proposition}[Charge-only repairs cannot remove the linear obstruction]
\label{prop:obstruction}
For the explicit inherited \(H_k\),
\[
 \lambda_{\max}(H_k)\ge C_k/4
       =\frac{k^{3/2}}{3\sqrt{2\pi}}+O(\sqrt k).
\]
For any real function \(f\) of total input charge and any \(t\ge0\),
positivity and normalization of
\[
 \widetilde\rho=\sigma^{1/2}[I-tH_k+f(Q)]\sigma^{1/2}
\]
necessarily imply \(tC_k\le4\).
\end{proposition}
\begin{proof}
Take the analytic input vector
\[
 \ket{\psi_k}=\left(\frac{\ket0+\ket1}{\sqrt2}\right)^{\otimes k}.
\]
All its colored wires are identical. For an oriented endpoint block,
the \(h>d\) term \(P(I-\SWAP)\) annihilates them, the \(h=d\)
term contributes \(1/2\), and the \(h<d\) term is zero.
The charge patterns have equal amplitudes. Including both orientations,
each pair contributes one quarter of its critical-height probability.
The inherited collision identity therefore gives exactly
\[
 \bra{\psi_k}H_k\ket{\psi_k}
 =\frac14\sum_{i+r\le k-2}\Prob(M_i=M_r)=C_k/4.
\]
Rayleigh's principle proves the spectral bound.

Normalization gives \(\Tr\sigma f(Q)=0\), since \(\Tr\sigma H=0\).
The charge \(Q\) has the same binomial distribution in \(\psi_k\) as
in \(\sigma\); hence \(\bra{\psi_k}f(Q)\ket{\psi_k}=0\).
Because \(\sigma\) is faithful, positivity requires
\(1-tC_k/4\ge0\).
\end{proof}

For the old coefficient \(t=(\ln2)b/2\), this requires
\(bC_k\le8/\ln2\), hence \(k=O(b^{-2/3})\) along any admissible
small-\(b\) sequence. For \(t=(\ln2)b/4\) the constant doubles but the
order is unchanged. Thus the new scale \(k\asymp1/[b\ln(1/b)]\) cannot
come from a sharper positivity estimate on this normalized linear family,
even with arbitrary charge-only diagonal repairs.
This does not determine the sharp positive spectral order or exclude
other input directions. The vector \(\psi_k\) is a spectral witness, not
a high-entropy or inward-drift communication input.

Normalization matters. The unnormalized choice
\(f(Q)=t\norm H I\) makes the bracket positive for arbitrary \(t\).
Normalizing that bracket changes the coefficient to
\(t/(1+t\norm H)\), so it does not evade the proposition.

\section{The weighted spectral scale is linear in temporal extent}

Let \(P_x\), \(x\in\{0,1\}^k\), be the projections onto complete
binary charge patterns and let \(\mathcal P(X)=\sum_xP_xXP_x\).
The color fibers inside these projections may have arbitrary dimension.

\begin{theorem}[Even-endpoint moments]\label{thm:moments}
Suppose \(J=\sum_{i<j}J_{ij}\) is Hermitian, each summand is Hermitian
with \(\norm{J_{ij}}\le2\), and
\[
 P_yJ_{ij}P_x=0
 \quad\hbox{unless}\quad y=x\mathbin{\mathrm{xor}}e_i
                                    \mathbin{\mathrm{xor}}e_j.
\]
Then, for every integer \(p\ge0\),
\begin{align}
 0\le\mathcal P(J^{2p})&\le(4p-1)!!\,k^{2p}I,\label{eq:moment}\\
 \mathcal P\!\left(e^{|J|/(4k)}\right)&\le\frac83I.\label{eq:pinchedexp}
\end{align}
Here \((-1)!!=1\). For every density \(\omega=\mathcal P(\omega)\),
\begin{align}
 \Tr[\omega\,1_{\{|J|>u\}}]&\le\frac83e^{-u/(4k)},\label{eq:tail}\\
 \Tr[\omega(|J|-u)_+]&\le\frac{32k}{3}e^{-u/(4k)}
 \qquad (u\ge0).\label{eq:integratedtail}
\end{align}
No product structure, commuting summands or tensor-local color action is required.
\end{theorem}
\begin{proof}
Expand \(J^{2p}\). A product can have a nonzero diagonal charge block
only if every endpoint appears an even number of times. There are \(2p\)
unordered edges in its ordered list. Orienting each edge yields \(2^{2p}\)
distinct lists of \(4p\) endpoint positions.

Every even-multiplicity endpoint list has a matching of equal labels.
There are \((4p-1)!!\) perfect matchings of the positions and at most
\(k^{2p}\) assignments of labels to the matching pairs. This overcounts:
self-edges and repeated representations need not be excluded. Therefore
the number of possibly nonzero edge products is at most
\((4p-1)!!\,k^{2p}/2^{2p}\).
Each product has norm at most \(2^{2p}\).
The same bound holds on every diagonal block. Positivity of the pinching
of \(J^{2p}\) proves \eqref{eq:moment}.

For \(0\le s<1/(2k)\), the nonnegative even power series and
\[
 \frac{(4p-1)!!}{(2p)!}
 =\frac{\binom{4p}{2p}}{4^p}\le4^p
\]
give
\[
 \mathcal P(\cosh(sJ))\le\frac{I}{1-4s^2k^2}.
\]
Spectral calculus gives \(e^{s|J|}\le2\cosh(sJ)\).
Putting \(s=1/(4k)\) proves \eqref{eq:pinchedexp}.
Taking the trace against \(\omega\), then applying the scalar spectral
Markov inequality and integrating its tail proves
\eqref{eq:tail}--\eqref{eq:integratedtail}.
\end{proof}

Apply this theorem to \(J=H\) and the actual \(\sigma\). It controls
the weight of large positive eigenvalues even though
Proposition~\ref{prop:obstruction} proves such eigenvalues exist.
The parity count concerns the input operator. It is not a classical
noise model for the coherent interaction.

\subsection{Counterexamples to omitted premises}

Total-charge preservation, offdiagonality and summand norm bounds do not
replace the named-endpoint rule. On the first two qutrits put
\[
 T_0=\sum_{a=1,2}(\ket{0a}\bra{a0}+\ket{a0}\bra{0a})
       \otimes I_{\rm other\ ports},
\]
and set every labelled \(J_{ij}=2T_0\).
All terms are Hermitian, offdiagonal, commute with \(Q,\sigma\) and
have norm two. But \(J=2\binom{k}{2}T_0\) and
\(\Tr\sigma T_0^2=1/2\). At \(k=4\),
\[
 \Tr\sigma J^2=72>3k^2=48.
\]
The labels are not the actual flipped endpoints.

The reference pinching premise is separate. On binary qubits set
\[
 J=2\sum_{i<j}
   (\sigma_i^+\sigma_j^-+\sigma_j^+\sigma_i^-).
\]
The named-endpoint rule and norm-two bound hold, and total charge is
preserved. The uniform Dicke vector with \(q\) ones is an eigenvector
of eigenvalue \(2q(k-q)\), since each word has \(q(k-q)\) allowed exchanges.
For \(k=4,q=2\), its pure density commutes with total charge (and \(J\)),
but its second moment is \(64>48\).
It is not diagonal in complete binary charge patterns.
For \(k=8,q=4\), its exponential moment is \(e>8/3\).
Thus commutation with total charge is weaker than the reference condition.

\section{A positive-part entropy and drift interface}

The following reusable argument is finite dimensional and independent of
the particular tail estimate.

\begin{lemma}[Positive-part normalization]\label{lem:repair}
Let \(\sigma\) be a faithful density commuting with \(Q\), let \(B=B^*\) commute with \(\sigma,Q\),
and suppose \(\Tr\sigma B=1\). Put
\[
 Z=B-I,\quad N=B_-,\quad \epsilon=\Tr\sigma N,\quad
 f=B_+/(1+\epsilon),\quad \rho=\sigma f.
\]
Then \(\rho\) is a normalized positive charge-preserving density,
and
\begin{align}
 \norm{\rho-\sigma B}_1&=2\epsilon,\label{eq:tracerepair}\\
 D_2(\rho\Vert\sigma)&\le\frac{\Tr(\sigma Z^2)}{\ln2}.\label{eq:chisquare}
\end{align}
For the actual \(-\log_2\sigma=kI+Q\), set
\(M=\Tr(\sigma ZA)\), \(V=\Tr(\sigma Z^2)\).
If \(M\le k/2\), then
\begin{equation}\label{eq:entropygeneric}
 S(\rho)-3k/2\ge M-k\epsilon-V/\ln2.
\end{equation}
If the exact bulk drift of the signed comparison satisfies
\(D_b(\sigma B)\le-\kappa_0 I\), then
\begin{equation}\label{eq:driftgeneric}
 D_b(\rho)\le
       \frac{-\kappa_0+k\epsilon}{1+\epsilon}I
\end{equation}
on that same bulk.
\end{lemma}
\begin{proof}
Joint spectral calculus makes \(\sigma B_\pm\) positive, with orthogonal
supports. Since \(\Tr\sigma B_+=1+\epsilon\), normalization follows and
the difference in \eqref{eq:tracerepair} is the sum of two opposite-sign
orthogonal terms, each of trace norm \(\epsilon\).

Write expectation with respect to \(\sigma\) as \(E\). Then \(Ef=1\) and
\[
 E(f-1)^2=Ef^2-1
   =\frac{EB_+^2}{(1+\epsilon)^2}-1
   \le EB^2-1=EZ^2.
\]
For \(u\ge0\), with the continuous convention at zero,
\[
 u\ln u-u+1\le(u-1)^2
\]
by \(\ln u\le u-1\). Since \([\rho,\sigma]=0\), taking expectations
proves \eqref{eq:chisquare}. This global bound does not require
\(\norm Z<1\) or a Taylor expansion.

The charge excess is
\[
 \Tr(\rho A)=\frac{M+\Tr(\sigma AN)}{1+\epsilon}
       \ge\frac{M-k\epsilon/2}{1+\epsilon}\ge M-k\epsilon.
\]
The final inequality uses \(M\le k/2\).
The exact identity
\(S(\rho)-3k/2=\Tr(\rho A)-D_2(\rho\Vert\sigma)\)
gives \eqref{eq:entropygeneric}.

For the drift use the exact relation
\[
 D_b(\rho)=\frac{D_b(\sigma B)+D_b(\sigma N)}{1+\epsilon}.
\]
The positive correction has trace \(\epsilon\), so
\eqref{eq:driftnorm} bounds \(D_b(\sigma N)\le k\epsilon I\).
This proves \eqref{eq:driftgeneric} on the entire bath-prefix operator
space, without choosing a prefix density.
\end{proof}

Noncommuting \(\sigma,B\) are outside this simple interface:
\(\sigma B_+\) need not be Hermitian and cannot be silently treated as a density.
The actual construction commutes because \(\sigma\) is a function of \(Q\).
Also \([\sigma,Q]=0\) is needed for the charge-preserving conclusion:
with \(B=I\), a faithful \(\sigma\) not commuting with \(Q\) is unchanged by
the repair and is not charge preserving.
The signed \(\sigma B\) is only an analytic comparison.
The endpoint response theorem is applied to this linear comparison, never
to the generally many-endpoint coherences in \(B_+\).

\section{Uniform control of the actual clipped input}

First fix any \(k\ge2\), \(0<b\le1/256\), with
\begin{equation}\label{eq:budget}
 r=bk\le1/64,\qquad
 \epsilon\le\frac{\ell vb^2}{256k}.
\end{equation}
We will then prove that \eqref{eq:newk} meets these conditions.
For the explicit \(B\) in \eqref{eq:newrho},
\[
 \Tr\sigma B=1,\qquad
 M=mb^2=\frac{3\ell vb^2}{32},\qquad
 V=\frac{\ell^2vb^2}{16}+\frac{9\ell^2v^2b^4}{256k}.
\]
The cross term vanishes since \(H\) is input-offdiagonal.
Here \(\Tr\sigma A^2=k/4\) and \(\Tr\sigma QH=0\).
The diagonal correction has norm
\begin{equation}\label{eq:diagonalcorrection}
 \eta_0=2mb^2=\frac{3\ell vb^2}{16}
       \le\frac{3\ell r^2}{128}<\frac12.
\end{equation}

\subsection{The tail pays for positivity}

Because \(H,A\) commute, joint spectral calculus and
\eqref{eq:diagonalcorrection} give
\[
 B_-\le\gamma b\left(H-\frac1{2\gamma b}\right)_+.
\]
The positive spectral tail is bounded by the two-sided tail in
Theorem~\ref{thm:moments}. Thus
\begin{equation}\label{eq:epsilonbound}
 \epsilon\le\frac{32\gamma}{3}bk
          \exp\!\left[-\frac1{8\gamma bk}\right]
 =\frac{8\ell}{3}r
          \exp\!\left[-\frac1{2\ell r}\right].
\end{equation}
This estimate needs small \(bk\), not small \(bk^2\).

\subsection{Entropy}

Write \(\mathcal B=\ell vb^2\) and \(u=vb^2/k\).
The bounds \(v\le k^2/8\), \(r\le1/64\), \(k\ge2\) give
\[
 u\le\frac{r^2}{8k}\le\frac1{65536},\qquad M\le k/2.
\]
By \eqref{eq:budget}, the charge excess is at least
\[
 M-k\epsilon\ge \frac{23}{256}\mathcal B.
\]
Also
\[
 D_2(\rho\Vert\sigma)\le V/\ell
 =\frac{\mathcal B}{256}(16+9u)
 \le\frac{17}{256}\mathcal B.
\]
Hence
\[
 S(\rho)-3k/2\ge\frac{3}{128}\mathcal B
                  \ge\frac{\ell vb^2}{64}.
\]
The weaker final constant is retained uniformly in the stated theorem.

\subsection{Exact operator drift}

The diagonal identity and all-order response in \eqref{eq:newresponse}
give the exact signed-comparison formula
\[
 D_b(\sigma B)=
 \gamma b^2(1-b)\left[\frac{3v}{4}I-F_k(b)\right].
\]
Because \(8b(k-2)\le1/8\), we have \(F_k(b)\ge7vI/8\), and therefore
\[
 D_b(\sigma B)\le-\kappa_0 I,\qquad
 \kappa_0=\frac{\ell vb^2(1-b)}{32}.
\]
The repair budget gives \(k\epsilon\le\ell vb^2/256\le\kappa_0/4\),
using \(1-b\ge1/2\); also \(\epsilon<1\).
Lemma~\ref{lem:repair} then gives
\[
 D_b(\rho)\le-\frac{3\kappa_0}{4(1+\epsilon)}I
          \le-\frac{\kappa_0}{4}I
          =-\frac{\ell vb^2(1-b)}{128}I.
\]
This proves \eqref{eq:exactdrift}. The correction was paid using the
full physical displacement bound. No pair-support assumption was imposed
on the clipped density and no nonlinear coefficient was dropped.

\subsection{Choosing the block before the coding horizon}

Let \(0<b\le1/256\), \(T=\ln(e/b)\), and choose \eqref{eq:newk}.
The function \(b\ln(e/b)\) is increasing on \((0,1)\). At \(b=1/256\),
\(T=1+8\ln2<7\), so \(16bT<7/16<1/2\).
Consequently the unfloored length exceeds two and
\[
 k\ge2,\qquad k\ge\frac1{32bT}.
\]
Since \(T>5\), \(r=bk\le1/(16T)<1/64\).
Equation~\eqref{eq:epsilonbound} gives
\[
 \exp[-1/(2\ell r)]
 \le\exp[-8T/\ell]\le\exp[-8T]\le b^8
 \le\frac{3b}{65536}.
\]
The last inequality follows from \(b^7\le2^{-56}<3/65536\).
Using \(v\ge k^2/32\), we obtain
\[
 \epsilon\le\frac{\ell b^2 k}{8192}
             \le\frac{\ell vb^2}{256k}.
\]
This proves every condition in \eqref{eq:budget} with the actual floor
and uniformly over the claimed parameter range.

\section{The original root and arbitrary-reference coding operation}

The density \(\rho\) commutes with \(Q\); therefore its actual bath channel
has homogeneous depth increments in \(\{-k,\ldots,k\}\), even at the
root. The exact bulk drift \eqref{eq:exactdrift} meets the inherited
operator-coding theorem, with \(0<\kappa\le k\).
That proof, reproduced below, retains the finite root remainder:
\[
 \Phi_\rho^*e^{tL_E}\le
 \lambda e^{tL_E}+e^{2kt}P_{<k},\qquad
 t=\frac{\kappa}{2k^2},\quad
 \lambda=1-\frac{\kappa^2}{4k^2}.
\]
No globally strict contraction is asserted.

The exact positive eigenvalues of \(\rho\) define the maximum-probability
spectral types. Their domination, logarithmic bath cutoff, joint compression
with its \(r^2\) positivity factor, auxiliary CPTP normalization and polar
repair give an isometric encoder and deterministic decoder.
The auxiliary normalization is not physically implemented.
For the inherited choices \(A_n=\lfloor\log_2(n+1)\rfloor\) and \(e_n\to0\),
the full-diamond error for any fixed normal \(\tau\) is bounded by
\[
 \min\{2,\ 4e_n+2\sqrt{1-\Tr(P_{\le A_n}\tau)}\}\longrightarrow0.
\]
In half-diamond convention this is exactly
\[
 \min\{1,\ 2e_n+\sqrt{1-\Tr(P_{\le A_n}\tau)}\}.
\]
The same code works pointwise for every fixed normal initializer; no
uniform finite-horizon supremum over all normal initializers is asserted.

At \(\tau_\pi\), the inherited exponential-moment/entropy bound makes
complementary information subextensive and proves the coherent-information
limit in Theorem~\ref{thm:newrate}. Finally,
\[
 \frac{S(\rho)-3k/2}{k}\ge\frac{\ell vb^2}{64k}
 \ge\frac{\ell b^2k}{2048}
 \ge\frac{\ell b}{65536\ln(e/b)}.
\]
This completes the new operational theorem. The bath remains one evolving
memory throughout all repeated blocks. Its analytic prefix minimizer is
never physically prepared.

\section{What changed, comparators, and the precise remaining boundary}

The preceding cubic floor for \(0<\beta\le1/8\) remains valid:
\[
 \Qcap\ge\frac32+\frac{\ln2}{8192}\beta^3,
\]
with its stronger unsimplified value
\((\ln2)v_{k_8}\beta^4/(32k_8)\), \(k_8=\lfloor1/(4\beta)\rfloor\).
The actual pair and better nonscalar four-call constructions, and the
restricted fixed-block quartic upper, are preserved unchanged.
Fixed-angle comparisons must take the maximum of all applicable retained
constructions; the simplified new global floor may lose constants relative
to an older unsimplified bound.

For example at the originally selected \(\beta=1/512\), the new block
has \(k=1215\) and \(v_k=181562.0497185\ldots\). Its unsimplified
certified excess is \(2.35512760108\ldots\times10^{-11}\), compared with
the preserved Cycle8 value \(4.76818050955\ldots\times10^{-12}\)
at \(k_8=128\). The ratio is \(4.93925848\ldots\).
These are evaluations of proved lower bounds, not measured capacities.
At \(\beta=1/16\) and \(1/64\), the old unsimplified bound is larger
than the new one; the source comparison retains the maximum.

The new achievable excess has order
\(\beta^2/\ln(1/\beta)\) as \(\beta\to0\). In particular it eventually
exceeds every fixed multiple of the cubic lower and is at least a
constant times \(\beta^{2+\delta}\) for every fixed \(\delta>0\), for
sufficiently small \(\beta\). This is an attainable lower, not an exact
capacity exponent, a quadratic lower without the logarithm, or a converse.
The inherited small-angle feedback upper and the question of equality
between assisted and unassisted capacities are unchanged.

\subsection{Strong ordinary comparisons and historical scope}

Even-multiplicity moment counting is a conventional probability mechanism.
Rudelson--Vershynin's Hanson--Wright theorem \cite{RV}, Theorem 1.1,
controls scalar quadratic forms of independent subgaussian variables.
It supplies a relevant ordinary comparison for exponential quadratic-chaos
tails. Our operator is not presented as such a scalar quadratic form:
the endpoint blocks contain input-controlled nonlocal color permutations.
The explicit grading proof above establishes exactly the property needed,
without asserting a new general concentration principle or priority.

Giovannetti--Burgarth--Mancini \cite{GBM}, Definition 1 and Theorem 2,
give zero-error quantum codes for multi-use channels with subexponential
effective environment dimension. That remains the closest ordinary coding
mechanism. The present infinite bath needs selected-input drift, root control,
type domination, and the actual initializer/compression repairs to reach its
hypotheses approximately on a code subspace.
The new work changes the admissible input family and attainable temporal
scale; it does not introduce a new generic quantum coding theorem.
Internal mathematical review does not establish external acceptance,
historical originality, efficient implementation or physical performance.

\subsection{Resources and limitations}

The input dimension is \(3^k\); its rank is \(d=\rank B_+\le3^k\).
No efficient diagonalization, state compiler, spectral-type encoder or
decoder is supplied. The inherited type probability may be as small as
\((n+1)^{-d}\), and its explicit bath cutoff contains
\[
 t^{-1}\big[(d+16)\ln(n+1)+\ln(1+C)\big],
 \qquad C=e^{2kt}/(1-\lambda).
\]
These are existence constants, not a practical blocklength estimate.
Exact functional calculus defines the density; synthesis precision and
end-to-end compilation cost are unmeasured.
At \(\beta=1/512\), \(\log_{10}(3^{1215})=579.70232448\ldots\);
this is an input-dimension upper bound on rank, not a claim that clipping
retains full rank. The selected drift is
\(\kappa=1.43073456\ldots\times10^{-8}\), while the inherited
Lyapunov parameter is \(t=4.84592308\ldots\times10^{-15}\) and
\(1-\lambda=3.46661481\ldots\times10^{-23}\).
The theoretical type-cutoff coefficient obtained by substituting
\(d\le3^{1215}\) is correspondingly enormous.
The physical call count for \(n\) repeated blocks is \(nk\), and both
the logical code and recovery syndrome are counted as in the inherited
proof. Small replays and transport timings do not measure discovery cost.

\subsection{Failures and the next owned question}

The positive-spectrum witness rejects an \(O(k)\) extreme-spectrum shortcut
for this \(H\), and rules out extending the normalized affine family to the
new block scale by charge-only repairs. It does not solve the sharp
\(\lambda_{\max}(H_k)\) problem. The stronger operational construction makes
that extremum unnecessary for the present rate.

The clipped input has higher endpoint coherences. The old four-endpoint
counterexample, nonscalar higher response and spectator-deletion failure
remain valid. An exponential/Gibbs input still has no automatic
pair-response theorem. Here their whole additional response is paid by
\(k\epsilon\), with no hidden classicalization.

The selected next question is whether this logarithmic loss can be removed
for the same operator and physical operation. At \(bk\) fixed, the present
tail estimate alone does not make \(\epsilon\) small enough to pay
\(k\epsilon\) against the available inward drift. A stronger route must
control the \emph{actual signed operator response and charge cost} of the
spectral repair, or establish a real obstruction for this positive family.
A fixed-\(bk\) numerical prefix expectation will not suffice.
No further capacity or positivity claim is made at this checkpoint.

The prior scalar certificate and exact all-order response were actually
used: they certify the affine comparison's full inward drift. The new
endpoint moment bound and positive-part entropy lemma remove its global
positivity restriction; the unchanged root/coding operation turns the repair
into the stronger communication rate. The evaluator's optional anisotropic
result was not used. This is a receiving mathematical consequence, not a
causal claim about prompts, agents or recursive amplification.

\appendix
\section{Preserved endpoint response proof}\label{app:response}
The following appendices reproduce the mathematical response and coding
arguments from Cycle8. Their depth notation \(L\) means the \(L_E\) used
above; the affine input operator in this return is \(B\).
% Mathematical response source extracted from unchanged Cycle8 main proof.
\subsection{Inherited susceptibility and the exact target}

Put \(b=\beta^2\), \(a=1-b\), and
\[
 \sigma=\operatorname{diag}(1/2,1/4,1/4)^{\otimes k},
 \qquad-\log_2\sigma=kI+Q.
\]
The bounded physical displacement
\[
 G_\beta=Q-W_k^*QW_k,\qquad \norm{G_\beta}\le k
\]
equals \(W_k^*LW_k-L\) on the finite-coordinate core, so its definition
does not subtract possibly infinite depth expectations.
For a signed charge-preserving input \(X\), write
\[
 D_b(X)=\Tr_A[(X\otimes I)G_\beta].
\]
It commutes with \(L\). On initial depth at least \(k\), its input-charge
pinching has the inherited exact first-\(k\)-color representation
\[
 \mathcal G_k(b)=G_0+bG_1+R_b,\qquad
 G_0=(2Q-kI)\otimes I_F,\qquad F=(\mathbb C^2)^{\otimes k}.
\]
At most \(k\) initial colors are consumed; final untouched tails agree in
every contributing equal-charge pair. The representation is tensor identity
on the seed and remaining tail. It is a bulk identity, not a cutoff or a
replacement of the root.

For a prefix density \(\zeta\), define
\[
 H_\zeta=\Tr_F[(I\otimes\zeta)G_1]+2Q-kI,\qquad
 T_k(X)=\Tr_A[(X\otimes I)G_1].
\]
The unchanged Cycle7 theorem gives
\begin{align}
 v_k&=\min_{\zeta}\Tr(\sigma H_\zeta^2)
 =\frac{k(k-1)}8-\frac18
   \sum_{n=0}^{k-2}(k-1-n)\Prob(S_n=0),\label{eq:vk}\\
 H&=H_{\zeta_{\rm sym}},\qquad
 \zeta_{\rm sym}=\frac{P_{{\rm sym},k}}{k+1},
 \qquad T_k(\sigma H)=v_kI_F.\label{eq:scalar}
\end{align}
The walk \(S_n\) is simple symmetric and starts at zero. The minimum admits
all coherent and entangled prefix densities. The analytic density
\(\zeta_{\rm sym}\) is never physically prepared. Useful inherited bounds are
\begin{equation}\label{eq:inherited}
 \frac{k^2}{32}\le v_k\le\frac{k(k-1)}8,\quad
 v_k\ge\frac N8,\quad N=\binom{k}{2},\quad
 \norm H\le k(k-1),\quad
 \Tr\sigma H=\Tr\sigma QH=0.
\end{equation}
Also \(H=H^*\), \([H,Q]=[H,\sigma]=0\), and \(H\) is input-offdiagonal.

We need its actual block structure, not just its norm. Write
\(H=\sum_{i<j}H_{ij}\), where \(H_{ij}\) connects complete binary charge
patterns differing by \(01\leftrightarrow10\) at \(i,j\).
On an oriented block, the Cycle7 normal form is
\begin{equation}\label{eq:Hblocks}
 K=
 \begin{cases}
 P(I-\SWAP),&h>d,\\
 P/2,&h=d,\\
 0,&h<d.
 \end{cases}
\end{equation}
Here \(h\) is the unmatched earlier push count, \(d\) the unmatched
middle pop count, and \(P\) a color-wire permutation. The reversed block
is \(K^*\). The two swapped wires in the first case are distinct input
color qubits. These are inherited bulk coefficient identities.

The preceding sufficient window \(bk^4\le2^{-12}\) came from bounding
\(\norm{R_b}\le4k^5b^2\) before contraction. We instead control the exact
quantity \(\Tr_A[(\sigma H\otimes I)\mathcal G_k(b)]\).

\subsection{Exact response from paired coherent histories}

\begin{theorem}[All-order endpoint response]\label{thm:response}
For every \(k\ge2\) and \(0\le b\le1\), there is a Hermitian operator
polynomial \(F_k(b)\) on the entire prefix \(F\) such that
\begin{align}
 \Tr_A[(\sigma H\otimes I)\mathcal G_k(b)]
   &=b(1-b)F_k(b),\qquad F_k(0)=v_kI,\label{eq:Fidentity}\\
 \norm{F_k(b)-v_kI}
   &\le N[1-(1-b)^{k-2}]
    \le N(k-2)b.\label{eq:Fbound}
\end{align}
In particular
\begin{equation}\label{eq:contracted}
 \norm{\Tr_A[(\sigma H\otimes I)R_b]}
 \le b^2v_k+b(1-b)N[1-(1-b)^{k-2}]
 \le b^2[v_k+N(k-2)].
\end{equation}
Every inequality is an operator norm bound, not an expectation at one
selected prefix.
\end{theorem}

\subsubsection{Charge selection and the cancellation of output charge}

For a stay mask \(M\subseteq\{0,\ldots,k-1\}\), let \(U_M\) be the ordered
word with \(I\) at \(M\) and \(U\) elsewhere. Its coherent amplitude is
\[
 \alpha^{k-|M|}(i\beta)^{|M|}U_M.
\]
On a bulk computational charge pattern \(x\), the output charge at a site
is \(x_\ell\) for a stay and \(1-x_\ell\) for a move. Suppose two initial
patterns differ only at endpoints \(i,j\), with opposite charges there.
A common output charge word forces the two histories to make identical
choices at every spectator and complementary choices at both endpoints.

Fix the common spectator stay set \(S\). Its paired scalar weight is
\[
 w_S(b)=b^{|S|}a^{k-2-|S|},\qquad
 \sum_{S\subseteq[k]\setminus\{i,j\}}w_S(b)=1.
\]
The endpoints supply \(ba\). Use a ket stay mask
\(s=(s_i,s_j)\), with ket input charges \(01\), and let \(R_s\) denote
the compressed raw cross-history operator. The bra mask is complementary:
\[
\begin{array}{c|c|c|c}
s&\text{bra stay mask}&\text{phase sign}&
       \text{output endpoint charge}\\ \hline
00&11&-1&1\\
01&10&+1&2\\
10&01&+1&0\\
11&00&-1&1
\end{array}
\]
Each \(R_s\) has norm at most one. This retains all input color and final
bath overlaps.

Keep the spectator calls fixed to \(U\) or \(I\), with only the endpoints
equal to \(V_\beta\). This auxiliary word is unitary. Its offdiagonal
input block in \(W^*W\) vanishes, and the four terms have a common \(ba\)
factor. Hence, for \(0<b<1\), and then as a polynomial identity,
\begin{equation}\label{eq:unitarity}
 -R_{00}+R_{01}+R_{10}-R_{11}=0.
\end{equation}
The spectator output charge is constant on the selected charge-pattern
block. Equation~\eqref{eq:unitarity} cancels that entire contribution to
\(W^*QW\), and also cancels the constant one in the endpoint charge.
The surviving block in \(Q-W^*QW\) is
\[
 R_{10}-R_{01},
\]
a difference of two contractions. There is no factor \(k\) from
\(\norm Q\).

Let \(B_{ij}^S\) be this block, its adjoint and their direct sum over
the complete charge patterns. The pattern pairs are disjoint, so
\begin{equation}\label{eq:Bnorm}
 \norm{B_{ij}^S}\le2.
\end{equation}
Equivalently it is the \(i,j\) endpoint component of the first coefficient
of the ordered word with spectator calls \(S\) skipped, embedded on the
original input space. This is an analytic representation only; the actual
communication process still uses and counts all \(k\) calls.

Define
\begin{equation}\label{eq:A}
 A_{ij,S}=\Tr_A[(\sigma H_{ij}\otimes I)B_{ij}^S].
\end{equation}
The original \(k\)-call \(H_{ij}\) remains in this formula. It is not replaced
by an optimizer for a shorter word. Exact history pairing now gives
\begin{equation}\label{eq:Bernstein}
 F_k(b)=\sum_{i<j}\ \sum_{S\subseteq[k]\setminus\{i,j\}}
          b^{|S|}a^{k-2-|S|}A_{ij,S}.
\end{equation}
Each \(A_{ij,S}\) is Hermitian by input partial-trace cyclicity. No
commutation of \(H_{ij}\) with \(B_{ij}^S\) is required.

The positive spectator weights arise after matching the two coherent
histories. Equation~\eqref{eq:Bernstein} does not make the physical
device a classical independent-error channel.

\subsubsection{The weighted trace norm that makes the estimate uniform}

\begin{lemma}[Endpoint trace budget]\label{lem:trace}
For the actual inherited \(H\),
\[
 \norm{\sigma H_{ij}}_1\le\frac14,\qquad
 \norm{A_{ij,S}}\le\frac12.
\]
\end{lemma}
\begin{proof}
Conditional on a complete charge pattern with \(q\) colors, the input
color weight is \(I/2^q\). In the first case of \eqref{eq:Hblocks},
\[
 |P(I-\SWAP)|=I-\SWAP,\qquad
 2^{-q}\Tr(I-\SWAP)=\frac12.
\]
The swap has normalized trace \(1/2\) on two qubits. In the critical
case \(|P/2|=I/2\), with the same normalized trace. The final case is
zero. The two orientations have identical singular-value sums, while
their combined endpoint-charge probability under \(\sigma\) is \(1/2\).
Since \(H_{ij}\) commutes with \(\sigma\),
\[
 \norm{\sigma H_{ij}}_1
 =\Tr(\sigma|H_{ij}|)\le\frac12\frac12=\frac14.
\]

For a Hermitian input operator \(X\) and bounded joint operator \(B\),
\[
 \norm{\Tr_A[(X\otimes I)B]}\le\norm X_1\norm B.
\]
Diagonalize \(X\), bound the norm of each corresponding input diagonal
block of \(B\), and sum with the absolute eigenvalues. This remains valid
with an infinite bath. Apply it to \eqref{eq:A} and \eqref{eq:Bnorm}.
\end{proof}

The empty-mask certificate is
\(\sum_{i<j}A_{ij,\varnothing}=v_kI\), by \eqref{eq:scalar}.
For each pair, the nonempty-mask mass is \(1-a^{k-2}\).
Lemma~\ref{lem:trace} and the triangle inequality give
\[
 \left\|\sum_Sw_S A_{ij,S}-A_{ij,\varnothing}\right\|
 \le1-a^{k-2}.
\]
Sum over \(N\) pairs and use Bernoulli's inequality to prove
\eqref{eq:Fbound}. Finally
\[
 \Tr_A[(\sigma H\otimes I)R_b]
 =ba(F_k(b)-v_kI)-b^2v_kI,
\]
because \(H\) is input-offdiagonal. This proves
Theorem~\ref{thm:response}.

\subsubsection{Reusable interface and its actual premises}

For any Hermitian pair-supported \(J=\sum J_{ij}\) commuting with
\(\sigma\), the same proof gives \eqref{eq:Bernstein} with \(J\) in place
of \(H\), and
\[
 \norm{F_J(b)-F_J(0)}
 \le4\!\left(\sum_{i<j}\norm{\sigma J_{ij}}_1\right)
      [1-a^{k-2}].
\]
A known nonscalar first operator certificate is allowed. The scalar
certificate and quarter trace budgets above belong to the actual Cycle7
optimizer. General total-charge preservation does not imply pair support,
and the root does not satisfy the bulk charge-flip premise. Concrete
counterexamples are preserved in the unchanged Cycle8 source \cite{Cycle8}.


\section{Inherited operator drift and the full coding consequence}\label{sec:inheritedcoding}

This section reproduces the actual Cycle5 proof \cite{Cycle5}, carried unchanged through Cycle6 apart from labels. Its hypotheses are verified by \eqref{eq:exactdrift}; no fresh-channel replacement is made.

\begin{lemma}[Homogeneous increments]\label{lem:grade}
For any charge-preserving $k$-call density $\rho$, its normal bath channel
has Kraus operators of homogeneous depth degree
$\Delta\in\{-k,\ldots,k\}$. Consequently there are positive effects
$A_\Delta$ such that
\[
 \sum_\Delta A_\Delta=I,\quad [A_\Delta,L]=0,\quad
 D_s(\rho)=\sum_\Delta\Delta A_\Delta,
 \tag{C1}
\]
and, as a positive quadratic-form identity,
\[
 \Phi_\rho^*e^{tL}
 =\left(\sum_\Delta e^{t\Delta}A_\Delta\right)e^{tL}.
 \tag{C2}
\]
\end{lemma}
\begin{proof}
Diagonalize $\rho$ in joint eigenvectors of $Q_k$, and take computational
output words $o$. Charge conservation gives each corresponding bath Kraus
operator degree $\Delta=q-N(o)$:
$[L,K]=\Delta K$. The exchange $R$ has degree zero and satisfies this
identity at the vacuum as well. Set $A_\Delta=\sum_{\deg K=\Delta}K^*K$.
Trace preservation proves their sum is $I$; homogeneity proves they commute
with $L$ and yields (C1)--(C2) on finite-coordinate vectors. Positive
quadratic forms and monotone limits extend the exponential identity.
\end{proof}

\begin{theorem}[Uniform operator drift gives actual reliable codes]
\label{thm:operatorcode}
Fix $k$, $0<\beta<1$, and a charge-preserving density $\rho$.
Suppose its exact drift obeys
\[
 P_{\ge k}D_s(\rho)P_{\ge k}\le-\kappa P_{\ge k},
 \qquad0<\kappa\le k.
 \tag{C3}
\]
Then one deterministic unassisted code sequence, independent of the bath
initializer, has rate at least $S(\rho)/k$ and pointwise vanishing
full-diamond error for every fixed normal initializer.
For the actual $\tau_\pi$, the product-block coherent-information limit
is also $S(\rho)/k$. No scalar depth chain or radial state preservation
is required.
\end{theorem}
\begin{proof}
Set
\[
 t=\kappa/(2k^2),\quad
 \lambda=1-\kappa^2/(4k^2),\quad
 C=\frac{e^{2kt}}{1-\lambda}.
 \tag{C4}
\]
Since $kt\le1/2$, for $|\Delta|\le k$ Taylor's inequality gives
$e^{t\Delta}\le1+t\Delta+k^2t^2$.
Multiply by each positive effect in (C1) and sum. On the high-depth sector
this is at most $\lambda I$. The low-depth sector reaches depth at most
$2k-1$, so (C2) gives the global form inequality
\[
 \Phi_\rho^*e^{tL}
 \le\lambda e^{tL}+e^{2kt}P_{<k}.
 \tag{C5}
\]
There are no low/high cross terms. Iteration yields
\[
 M_m\le\lambda^m M_0+C(1-\lambda^m).
 \tag{C6}
\]
For any coherent initializer supported in $P_{\le A}$,
\[
 \Pr(L_m>H)\le(1+C)e^{-t(H-A)}.
 \tag{C7}
\]

Let $d=\rank\rho\le3^k$, using only positive eigenvalues. A
maximum-probability exact spectral type $T_m$ of $\rho^{\otimes m}$
has probability $w_m\ge(m+1)^{-d}$, rank $D_m$, and
\[
 T_m/D_m\le\rho^{\otimes m}/w_m,\qquad
 \log_2D_m/m\longrightarrow S(\rho).
 \tag{C8}
\]
For completeness, with type counts $n_i$, maximality against a transfer
from count $i$ to $j$ gives
$n_i/p_i\le(n_j+1)/p_j$. Summing against $p_i$ shows
$n_j\ge mp_j-1$, hence every frequency tends to $p_j$.
The type scalar eigenvalue and the polynomial mass give the rank limit.
This is an operator statement in the actual eigenbasis, even if eigenvectors
are entangled across calls.

Choose
\[
 A_m=\lfloor\log_2(m+1)\rfloor,\quad
 H_m=A_m+\left\lceil
 \frac{(d+16)\ln(m+1)+\ln(1+C)}t\right\rceil,\quad
 \epsilon_m=(m+1)^{-16}.
 \tag{C9}
\]
Apply (C7) to the uniform state on $P_{\le A_m}$, and then (C8).
Mean overflow for the type subspace and that initial bath space is at
most $\epsilon_m$. The inherited joint compression lemma, reproduced
in Appendix~\ref{app:compression}, gives
\[
 K_m\ge\left\lceil
 \frac{\lfloor D_m/2\rfloor}{(r_{A_m}r_{H_m})^4}
 \right\rceil,\quad
 e_m=r_{A_m}\sqrt{2\epsilon_m}+(1-2\epsilon_m)^{-1/2}-1\to0.
 \tag{C10}
\]
It retains spectral expurgation, the $r^2$ positive-block factor,
the auxiliary CPTP normalization and its polar repair.
The actual encoder is isometric; no auxiliary normalization is implemented.
For every fixed normal $\tau$, the full-diamond error is at most
\[
 \min\{2,\ 4e_m+2\sqrt{1-\Tr(P_{\le A_m}\tau)}\}\longrightarrow0.
 \tag{C11}
\]
In half-diamond convention this is
$\min\{1,2e_m+\sqrt{1-\Tr(P_{\le A_m}\tau)}\}$.
This follows by comparing $\tau$ outside the type domination with its
normalized coordinate truncation. Thus no type-probability factor multiplies
an uncontrolled initial tail. Since $\log r_{A_m}r_{H_m}=O(\log m)$
for this fixed device and block, (C10) proves the rate.

For $\tau_\pi$, $t\le1/2<\ln3$ ensures a finite initial exponential
moment. Equations (C6)--(C7) bound mean depth uniformly. The bath Gibbs
density $g=4^{-L}/4$ has trace one and
$-\log_2g=2L+2$, so $S(E_m)\le2\E L_m+2$.
This entropy-energy inequality also follows by finite cutoffs and monotone
limits; finite mean guarantees finite entropy. The initial bath entropy
is finite. Unitary entropy conservation, subadditivity and Araki--Lieb give
\[
 mS(\rho)-2S(E_m)\le\Ic(\rho^{\otimes m},\mathcal N_{km}^{\tau_\pi})
 \le mS(\rho).
 \tag{C12}
\]
Divide by $km$. Extra data-independent calls may fill incomplete final
blocks; all count and do not alter already emitted output marginals.
\end{proof}

The initializer-independent code quantifier is pointwise in the normal initializer.
It is distinct from the inherited uniform all-normal restricted upper bound.
No invariant initializer is prepared. The boundary term in (C5) is essential:
a global strict exponential contraction would contradict trace preservation
and $e^{tL}\ge I$.


\section{The inherited joint compression operation}\label{app:compression}

This appendix restates the actual inherited proof \cite{Cycle2}, rather than
using a summary as a coding theorem. It records the normalization and coherent
initializer quantifiers needed by the new receiving use. The underlying
finite-error mechanism is established \cite{KLV,GBM,Tverberg}.

\begin{lemma}[Finite-error-set code]\label{lem:finitecode}
A CPTP map from a $D$-dimensional input space to a finite output space, with
at most $k$ Kraus operators, has an exact subspace code of dimension at least
$\lceil D/k^4\rceil$.
\end{lemma}
\begin{proof}
The self-adjoint operator space spanned by $A_i^*A_j$ contains $I$ and has
complex dimension $d\le k^2$. Choose a Hermitian basis $I,H_2,\ldots,H_d$.
Greedily choose orthonormal vectors $v_1,\ldots,v_N$ with all cross matrix
elements of each $H_j$ zero: after $m$ choices avoid the span of at most
$md$ vectors $H_jv_a$. This yields $N\ge\lceil D/d\rceil$.
Their diagonal expectation vectors lie in $\mathbb R^{d-1}$. The affine
Tverberg theorem partitions any $(r-1)d+1$ such points into $r$ groups with
intersecting convex hulls. Choose $r=\lceil N/d\rceil$. Square roots of the
corresponding convex weights form orthonormal code vectors with identical
diagonal expectations and zero cross expectations. Thus
$C^*A_i^*A_jC=\gamma_{ij}I_r$ and $r\ge\lceil D/k^4\rceil$.

Choose syndrome vectors of Gram matrix $\gamma$ and send
$A_iC\ket\psi\mapsto\ket\psi\otimes\ket{\xi_i}$ on the error-image span.
The Gram identities make this map isometric; extend it with an external
syndrome space to a recovery isometry. All outcomes are recovered. Nothing
requires equal input and output dimensions.
\end{proof}

\begin{lemma}[Joint compression from averaged overflow]\label{lem:compression}
Let $W:A\otimes E_0\to B\otimes E$ be an isometry, with
$\dim A=D$, $\dim E_0=r<\infty$, and let $P$ be a rank-$d$ bath projection.
Define
\[
 F=W^*(I_B\otimes(I-P))W,\qquad L_A=\frac1r\Tr_{E_0}F.
\]
If $D^{-1}\Tr L_A\le\epsilon$, $0<\epsilon<1/2$, there is an isometric
encoder and deterministic decoder of dimension
\[
 K\ge\left\lceil\frac{\lfloor D/2\rfloor}{(rd)^4}\right\rceil
\]
whose recovered logical channel has full-diamond error at most $4e$ for
every initializer supported in $E_0$, provided $e<1$, where
\[
 e=r\sqrt{2\epsilon}+(1-2\epsilon)^{-1/2}-1.
\]
The same code is valid for coherent initializers and arbitrary external
references. More strongly its recovered joint isometry is within $2e$ of
$I_{\rm logical}\otimes Z_0$ for a bath/syndrome isometry $Z_0$.
\end{lemma}
\begin{proof}
Let $Q=1_{[0,2\epsilon]}(L_A)$. Eigenvalue counting gives
$\dim QA\ge D/2$ and $QL_AQ\le2\epsilon Q$.
For any positive operator $X$ on $A\otimes\mathbb C^r$,
\[
 X\le r(\Tr_{\mathbb C^r}X)\otimes I_r.
 \tag{A1}
\]
Indeed pinching into $r$ diagonal bath blocks gives
$X\le r\sum_v(I\otimes|v\rangle\langle v|)X(I\otimes|v\rangle\langle v|)$
by Cauchy--Schwarz, and each positive diagonal block is at most their sum.
For $F_Q=(Q\otimes I)F(Q\otimes I)$, (A1) implies
$F_Q\le r^2 QL_AQ\otimes I\le2r^2\epsilon I$. The discarded isometry
amplitude is therefore at most $r\sqrt{2\epsilon}$, uniformly on joint
inputs, including entangled references.

For bases $v$ of $E_0$ and $w$ of $PE$, put
\[
 B_{wv}=\frac1{\sqrt r}\bra w(I\otimes P)W\ket v\big|_{QA}.
\]
Then $\sum B_{wv}^*B_{wv}=I_{QA}-QL_AQ$. The auxiliary operator
$J=(I_{QA}-QL_AQ)^{-1/2}$ satisfies
$\norm{J-I}\le(1-2\epsilon)^{-1/2}-1$, and $B_{wv}J$ define a CPTP map
with at most $rd$ Kraus operators. Lemma~\ref{lem:finitecode} gives an
isometric code inclusion $C$ and recovery isometry $D_{\rm rec}$ with
$D_{\rm rec}B_{wv}JC\ket\psi=\ket\psi\otimes\ket{\xi_{wv}}$.
Consequently
\[
 (D_{\rm rec}\otimes I)(I\otimes P)W(JC\otimes I_{E_0})
 =I_{\rm logical}\otimes Z,\quad
 Z\ket v=\sqrt r\sum_w\ket{\xi_{wv}}\otimes\ket w.
 \tag{A2}
\]
This uses all initial and final bath indices, not one pure initial bath.

The actual encoder is $C$, never $JC$. Its recovered joint isometry
$H=(D_{\rm rec}\otimes I)W(C\otimes I)$ differs from $I\otimes Z$ by at
most $e$. Since $H$ is isometric, every singular value of $Z$ lies in
$[1-e,1+e]$. Its polar part $Z_0$ is therefore isometric and
$\norm{Z-Z_0}\le e$, giving $\norm{H-I\otimes Z_0}\le2e$.
Two isometries differing in norm by $u$ induce channels at full-diamond
distance at most $2u$: apply the pure-state trace-distance inequality after
tensoring a reference, then trace unused outputs. This proves $4e$, or $2e$
in half-diamond convention. The recovery syndrome is retained in the dilation;
discarding it to define a logical marginal is not an internal bath reset.
\end{proof}

The $r$ in (A1) cannot be removed: for the maximally entangled projector
$X=|\Omega_r\rangle\langle\Omega_r|$ one has $\norm X=1$ but
$(\Tr_E X)\otimes I=I/r$. Thus a bath-averaged effect of norm $1/r^2$ can
hide a joint failure of probability one. Likewise the projected channel is
not automatically trace preserving. The auxiliary $J$ and the polar repair
are mathematical comparisons; implementing $J$ as a deterministic encoder
or omitting its error would be invalid.


\begingroup\small
\begin{thebibliography}{9}
\setlength{\itemsep}{2pt}\setlength{\parskip}{0pt}
\bibitem{Cycle8} Astra Pro 1, \emph{Exact nonlinear response of temporal
quantum coherences}, source-complete Cycle8 checkpoint, 22 October 2026.
Exact original archive preserved unchanged in this return.
\bibitem{Cycle7} Astra Pro 1, \emph{Sharp temporal susceptibility for a
recurrent quantum channel}, source-complete Cycle7 checkpoint, 22 October 2026.
Exact native archive and earlier chain preserved inside Cycle8.
\bibitem{Cycle5} Astra Pro 1, \emph{A quartic information law for finite
temporal blocks}, source-complete Cycle5 checkpoint, 22 October 2026.
Original preserved within the native predecessor chain.
\bibitem{Cycle2} Astra Pro 1, original Cycle2--Cycle4 proof/source chain.
The actual joint-compression proof is reproduced in this return.
\bibitem{RV} M. Rudelson and R. Vershynin,
\emph{Hanson--Wright inequality and sub-gaussian concentration},
Electronic Communications in Probability \textbf{18}, no. 82 (2013),
Theorem 1.1, \url{https://arxiv.org/abs/1306.2872},
\url{https://doi.org/10.1214/ECP.v18-2865}.
\bibitem{GBM} V. Giovannetti, D. Burgarth and S. Mancini,
\emph{Communication Through a Quantum Link},
Phys. Rev. A \textbf{79}, 012311 (2009),
\url{https://arxiv.org/abs/0810.0170}, Definition 1 and Theorem 2.
\bibitem{KLV} E. Knill, R. Laflamme and L. Viola,
\emph{Theory of Quantum Error Correction for General Noise},
Phys. Rev. Lett. \textbf{84}, 2525 (2000),
\url{https://arxiv.org/abs/quant-ph/9908066}.
\bibitem{Tverberg} H. Tverberg, \emph{A generalization of Radon's theorem},
J. London Math. Soc. \textbf{41}, 123--128 (1966),
\url{https://doi.org/10.1112/jlms/s1-41.1.123}.
\end{thebibliography}
\endgroup
\end{document}
